% Part: computability
% Chapter: computability-theory
% Section: s-m-n

\documentclass[../../../include/open-logic-section]{subfiles}

\begin{document}

\olfileid{cmp}{thy}{smn}
\olsection{$s$-$m$-$n$ సిద్ధాంతం}

\begin{explain}
తరువాతి సిద్ధాంతాన్ని ``$s$-$m$-$n$ సిద్ధాంతం'' అంటారు; దానికి కారణం
కొద్దిసేపటిలో స్పష్టమవుతుంది. సిద్ధాంతం సరిగ్గా ఏమి చెబుతుందో అర్థం
చేసుకోవడమే కష్టమైన భాగం; ప్రకటన అర్థమైన తరువాత అది చాలావరకు స్పష్టంగా
అనిపిస్తుంది.
\end{explain}

\begin{thm}
\ollabel{thm:s-m-n}
  ప్రతి సహజ సంఖ్యల జత $n$,~$m$కు, ప్రతి అనుక్రమం
  $e$, $a_0$, \dots, $a_{m-1}$, $y_0$, \dots, $y_{n-1}$కు
  \[
  \cfind{s^m_n(e, a_0, \dots, a_{m-1})}[n](y_0, \dots, y_{n-1}) \simeq
  \cfind{e}[m+n](a_0, \dots, a_{m-1}, y_0, \dots, y_{n-1}).
\]
  అయ్యే ఆదిమ పునరావృత్త ప్రమేయం~$s^m_n$ ఒకటి ఉంటుంది.
\end{thm}

\begin{explain}
$s^m_n$ అనేది \emph{ప్రోగ్రాములపై} పనిచేస్తుందని భావించడం ఉపయుక్తం.
అంటే, ఒక $(m+n)$-స్థాన ప్రమేయానికి ప్రోగ్రాము~$e$తో పాటు స్థిర
ఇన్‌పుట్‌లు $a_0$, \dots, $a_{m-1}$లను $s^m_n$ తీసుకుని, మిగిలిన
ఆర్గుమెంట్ల $n$-స్థాన ప్రమేయానికి ప్రోగ్రాము
$s^m_n(e,a_0,\dots,a_{m-1})$ను తిరిగి ఇస్తుంది.
\iftag{TMs}{$e$ను ఒక ట్యూరింగ్ యంత్రపు వివరణగా భావిస్తే,
$s^m_n(e,a_0,\dots,a_{m-1})$ అనేది ఇన్‌పుట్ $y_0$, \dots,
$y_{n-1}$కు ముందు $a_0$, \dots, $a_{m-1}$లను చేర్చి, తరువాత
యంత్రం~$e$ను నడిపే ట్యూరింగ్ యంత్రం. అప్పుడు ప్రతి $s^m_n$ తగిన
ట్యూరింగ్ యంత్రానికి సంకేతాన్ని కనుగొనే ఆదిమ పునరావృత్త ప్రమేయమే.}{}
\sourcecorrection{OLTECOMTHY-001}{ఆంగ్ల మూలంలోని వివరణ సిద్ధాంతపు ప్రోగ్రాము సూచికను మొదట $e$గా పేర్కొన్నప్పటికీ, తిరిగి ఇచ్చే ప్రోగ్రామును $s^m_n(x,a_0,\dots,a_{m-1})$గా రాసింది. సిద్ధాంత ప్రకటనకూ అదే వివరణలోని ట్యూరింగ్ యంత్ర వాక్యానికీ అనుగుణంగా ఆ చరాన్ని $e$గా సరిచేశాం.}
\end{explain}

\end{document}
